\chapter{Összegzés}

\section{Mit mutat be a dolgozat?}

A dolgozat egy olyan rendszert mutat be, amely valós League of Legends-mérkőzésadatokból indul ki, majd ezekből játékosprofilokat, kapcsolatokat, gráfokat, klasztereket, elemzési eredményeket és végül egy szimulációs kísérletet épít. A rendszer a játékosok ölés- és halálszámain túl azt is rögzíti, hogy egy játékos kikkel kerül egy meccsbe, kikkel kerül szembe, milyen ismétlődő kapcsolatok rajzolódnak ki körülötte, és ezekből milyen nagyobb hálózati minták olvashatók ki.

Ha valamit le kellene szűkíteni a projekt lényegéből, az inkább az egész lánc megléte, mint bármely egyedi komponensé. Egy Dijkstra-implementáció vagy egy React-oldal önmagában egy szokásos projekt. A rendszer érdekessége az összekötöttség: a nyers meccsadatoktól a gráfon, klasztereken és Rust-analitikán át eljut egy reprodukálható szimulációs kísérletig, és minden lépés ellenőrizhető.

\section{A fő eredmények}

Az adatgyűjtési réteg két eltérő queue-típust kezel párhuzamosan. A \textit{flexset} és a \textit{soloq} külön kezelése azért fontos, mert a két queue nem ugyanazt a jelenséget méri: a Flex Queue alkalmasabb az ismétlődő kapcsolati minták vizsgálatára, a SoloQ pedig kontrollként és egyéni teljesítményprofilként értelmezhető.

Az \textit{opscore} és \textit{feedscore} metrikák szerepkör-érzékeny pontszámokat adnak. Nem tökéletes igazságmérők, de lehetővé teszik, hogy a teljesítménymutatók gráfanalitikai bemenetként is használhatók legyenek.

A \textit{flexset} gráf szerkezete core-periphery jellegű: vannak sűrűbb, erősebben összekapcsolt magrészek és sok lazább, kisebb periférikus csoport. Ez a queue-mechanikából következik.

Az asszortativitási elemzés az \textit{opscore} esetében pozitív kapcsolatot mutat (social-path: 0,235): az összekapcsolt játékosok hasonlóbb teljesítményprofilt hordoznak, mint amit véletlenszerű kapcsolatoknál várnánk.

A Rust-alapú algoritmikus réteg (útvonalkeresés, asszortativitás, betweenness centralitás) determinisztikus és újrafuttatható: ugyanabból a bemeneti adatból mindig ugyanaz az eredmény jön ki.

\section{A Tribal NeuroSim szerepe}

A Tribal NeuroSim átviteli kísérlet: azt vizsgálja, hogy a gráfból és a klaszterekből kinyert teljesítményprofilok átvihetők-e egy evolúciós, többügynökös szimuláció kezdeti paramétereibe.

A négy futtatás alapján a Flex Queue-ból származó profilok egységesebb viselkedést mutattak: mindkét vizsgált seed mellett hasonló, Warband/Combat jellegű győztes archetípus jelent meg. A SoloQ futtatások változékonyabbak voltak, eltérő seedek mellett eltérő győztes profilok alakultak ki. A flexset-priorok tehát egységesebb kimenetet adtak, a soloq-priorok változékonyabbat, azonos szimulációs szabályrendszer mellett.

\section{Visszatekintés: ami jól működött}

Visszatekintve a legjobb döntés az volt, hogy a rendszer nem erőltette mindent egyetlen technológiába. A Python gyors adatgyűjtésre és kísérletezésre maradt, a Rust ott vette át a szerepet, ahol teljesítmény és determinizmus kellett, a Node/Express és a React pedig bemutathatóvá tette az egészet. Ha mindezt egyetlen nyelvre kellett volna szűkíteni, valamelyik réteg szenvedett volna miatta.

A projekt erőssége az összefüggő lánc: a nyers API-adatoktól a normalizált játékosprofilokon, a gráfépítésen és a klaszterezésen át egészen az analitikai eredményekig minden lépés ellenőrizhető és újrafuttatható. Az asszortativitás, a betweenness centralitás és a core-periphery értelmezés nemcsak vizualizációt tesz lehetővé, hanem valódi kérdéseket is felvet a hálózatról.

\section{Korlátok és nyitott kérdések}

Az eredmények értelmezésekor néhány korlátot érdemes szem előtt tartani.

Az adatkészlet nem véletlenszerű mintavétel. A seed-játékosok kiválasztása, a premade-probe mechanizmus és a queue-típus együtt meghatározzák, milyen játékospopuláció kerül be a gráfba. A kapott szerkezet az adott gyűjtési stratégia terméke, nem általánosan érvényes EUNE Master+ leírás.

Az opscore és feedscore hasznos és értelmezhető mutatók, de nem külső igazságmérők. A szerepkör-elvárások és a görbe paraméterei a modell részei; más konstrukciós döntéssel más eredményt kapnánk.

A gráfban az ellenség-kapcsolat sem egyenlő valódi konfliktussal. Az ismételt szembekerülés jöhet a matchmaking MMR-sávjából, az azonos napszakból vagy abból, hogy a Master+ Flex Queue aktív játékospoolja eleve szűk. Ezért maradt az előjeles gráf kísérleti opciónak.

A temporal consistency és a community cohesion vizsgálatok kiestek a jelenlegi scope-ból. A kiesett vizsgálatok nem érdektelenek; az adatok és a módszer azonban nem tette lehetővé a védhető állításokat.

\section{Zárszó}

A projekt végeredménye egyszerűen megfogalmazható: kompetitív játékadatokból lehetséges olyan rendszert felépíteni, amely szoftvermérnöki szempontból is megállja a helyét és kutatói kérdéseket is tud feltenni. A PremadeGraph adatgyűjtési, adatmodellezési, gráfanalitikai és szimulációs lánca ezt igazolja.

A lépések közötti átjárhatóság nem volt előre tervezett. A nyers meccsadatokból normalizált játékosprofil lett, abból kapcsolati gráf, a gráfból klaszterek és analitikai mutatók, a klaszterekből szimulációs priorok; minden egyes lépés az előző kimenetét vette bemenetnek.

Ahogy az ABBA megfogalmazta évtizedekkel ezelőtt: \textit{the winner takes it all}.
A Tribal NeuroSimben ez nem metafora, hanem a mechanika közvetlen következménye.

\section{Ahol ez máshol is jön számításba}
\label{sec:osszegzes-transfer}

A módszertan nem kötődik a League of Legends-hez. A Riot API azért volt kényelmes választás, mert nyilvánosan elérhető, gazdag és strukturált. De ugyanez a pipeline ugyanúgy lefutna GitHub-kollaborációs adatokon, NBA-teljesítménystatisztikákon vagy PubMed-szerzői hálózatokon, bárhol, ahol ismétlődő interakció, mérhető teljesítmény és valamilyen csoportszerkezet van.

A gráfklaszterezés például ugyanolyan természetesen leír egy vállalati szervezeti hálózatot, mint egy játékospopulációt: a kérdés mindkét esetben az, hogy a formális struktúra mögött milyen tényleges munkacsoportok rajzolódnak ki. A betweenness centralitás az infrastruktúra-tervezésben vagy az epidemiológiában is ugyanazt méri: melyik csomópont kiesése szakítja szét leginkább a hálózatot.

A Tribal NeuroSim öt artifaktja sem LoL-specifikus fogalom. Combat, resource, map\_objective, risk, team: ezek más doménekben is lefordíthatók, és az evolúciós szimuláció keretrendszere bármilyen klaszterprofilból kiindulhat, ahol csoportszintű teljesítménymutatók kinyerhetők. A videojáték itt eszköz volt, nem korlát.
